\chapter{Small-Angle Limits and Trigonometric Derivatives}
\label{ch:ch17-calculus-connections}

\section{From Geometry to Analysis}

The Introduction to this book described trigonometry not merely as a
collection of formulas, but as a bridge between geometry and calculus.
The preceding chapters have gradually built that bridge. Chapter~15
showed how scaling and rotation combine into amplitwists; Chapter~16
showed that complex numbers provide an algebraic language for those same
transformations. Geometry became algebra. This chapter takes the next
step: algebra now becomes analysis. The central question is simple: how
do trigonometric functions behave under infinitesimal change? To answer
it requires understanding what happens as an angle becomes extremely
small. The resulting limit leads directly to the derivatives of the sine
and cosine functions and establishes the first substantial connection
between trigonometry and calculus.

\section{The Geometry of Small Angles}

Consider the values of $\sin\theta$, $\theta$, and $\tan\theta$ for
progressively smaller angles.

\begin{table}[h]
\centering
\begin{tabular}{c|c|c}
$\theta$ & $\sin\theta$ & $\tan\theta$\\
\hline
$0.5$ & $0.4794$ & $0.5463$\\
$0.1$ & $0.0998$ & $0.1003$\\
$0.05$ & $0.0500$ & $0.0500$\\
$0.01$ & $0.0100$ & $0.0100$
\end{tabular}
\caption{Values of $\sin\theta$ and $\tan\theta$ as $\theta \to 0$ (radians).}
\end{table}

As $\theta$ approaches zero, the three quantities appear to become
indistinguishable, $\sin\theta \approx \theta \approx \tan\theta$. This
observation suggests an important limit, but before proving it a crucial
remark is needed: the approximation works only when $\theta$ is measured
in \emph{radians}. If degrees are used instead, the numerical agreement
disappears, and this fact is not accidental. Chapter~4 introduced radians
because they arise naturally from arc length, and this chapter reveals a
second reason: radians are the angular unit that makes the calculus of
trigonometric functions take its simplest form.

\section{The Sector-Area Sandwich}

Consider a unit circle and an angle $\theta$ satisfying $0 < \theta <
\pi/2$. Three regions can be constructed at this angle: the triangle
formed by the radius and the chord, the circular sector subtended by the
angle, and the larger triangle formed by the radius and the tangent
line. The first of these regions lies entirely inside the sector, while
the sector in turn lies entirely inside the outer triangle, giving
\[
\text{Area}_{\text{inner}} < \text{Area}_{\text{sector}} <
\text{Area}_{\text{outer}}.
\]
Each area can now be computed directly. The inner triangle has base $1$
and height $\sin\theta$, so $\text{Area}_{\text{inner}} =
\tfrac12\sin\theta$. The sector area was derived in Chapter~5:
$\text{Area}_{\text{sector}} = \tfrac12\theta$. The outer triangle has
base $1$ and height $\tan\theta$, giving $\text{Area}_{\text{outer}} =
\tfrac12\tan\theta$. Combining these results yields
\[
\frac12\sin\theta < \frac12\theta < \frac12\tan\theta,
\]
and multiplying by $2$ gives the fundamental inequality
\[
\sin\theta < \theta < \tan\theta.
\]
This geometric sandwich is the key to everything that follows.

\section{The Fundamental Small-Angle Limit}

The inequality $\sin\theta < \theta < \tan\theta$ now determines the
behavior of $\sin\theta/\theta$ as $\theta$ approaches zero.

\begin{theorem}
$\displaystyle\lim_{\theta\to0} \frac{\sin\theta}{\theta} = 1$.
\end{theorem}

\begin{proof}
For $0 < \theta < \pi/2$, $\sin\theta < \theta < \tan\theta$. Dividing
through by $\sin\theta$ gives $1 < \theta/\sin\theta < 1/\cos\theta$, and
taking reciprocals reverses the inequalities:
\[
\cos\theta < \frac{\sin\theta}{\theta} < 1.
\]
As $\theta \to 0$, $\cos\theta \to 1$, and the right-hand bound is
already $1$; by the Squeeze Theorem, $\lim_{\theta\to0} \sin\theta/\theta
= 1$.
\end{proof}

This result is the most important limit in elementary trigonometry, and
the derivatives of sine and cosine ultimately depend on it.

\section{Corollary Limits}

Several additional limits follow immediately.

\begin{theorem}
$\displaystyle\lim_{\theta\to0} \frac{\tan\theta}{\theta} = 1$.
\end{theorem}

\begin{proof}
Since $\tan\theta = \sin\theta/\cos\theta$, $\tan\theta/\theta =
(\sin\theta/\theta)\cdot(1/\cos\theta)$. Taking limits, $\lim_{\theta\to0}
\tan\theta/\theta = 1\cdot1 = 1$.
\end{proof}

\begin{theorem}
$\displaystyle\lim_{\theta\to0} \frac{1-\cos\theta}{\theta} = 0$.
\end{theorem}

\begin{proof}
Using the identity $1-\cos\theta = 2\sin^2(\theta/2)$ from Chapter~10,
\[
\frac{1-\cos\theta}{\theta} = \sin\left(\frac{\theta}{2}\right)\cdot
\frac{2\sin(\theta/2)}{\theta}.
\]
As $\theta \to 0$, $\sin(\theta/2) \to 0$, while
$\sin(\theta/2)/(\theta/2) \to 1$ by the fundamental limit. The entire
product therefore approaches $0$.
\end{proof}

\section{The Derivative}

The limits established above describe the behavior of trigonometric
functions under extremely small changes. Calculus formalizes this idea
through the concept of a derivative. Suppose a function $f(x)$ is
defined on an interval, and let $x$ and $x+h$ be two nearby points; the
average rate of change of the function between them is
$(f(x+h)-f(x))/h$, the slope of the secant line joining $(x,f(x))$ and
$(x+h,f(x+h))$. As $h$ becomes smaller, the secant line approaches the
tangent line at $x$.

\begin{definition}
The \emph{derivative} of a function $f$ at $x$ is
\[
f'(x) = \lim_{h\to0} \frac{f(x+h)-f(x)}{h},
\]
provided the limit exists.
\end{definition}

The derivative measures the instantaneous rate of change of the
function, and gives the slope of its tangent line. The goal of the next
two sections is to compute the derivatives of the sine and cosine
functions.

\section{The Derivative of Sine}

Begin with $f(x) = \sin x$. By the definition of the derivative,
\[
f'(x) = \lim_{h\to0} \frac{\sin(x+h)-\sin x}{h}.
\]
The sum formula from Chapter~8 gives $\sin(x+h) = \sin x\cos h + \cos
x\sin h$, and substituting into the difference quotient,
\[
f'(x) = \lim_{h\to0} \frac{\sin x\cos h + \cos x\sin h - \sin x}{h}.
\]
Grouping terms,
\[
f'(x) = \lim_{h\to0} \left[\sin x\,\frac{\cos h-1}{h} + \cos
x\,\frac{\sin h}{h}\right].
\]
Since $\sin x$ and $\cos x$ do not depend on $h$, they may be treated as
constants:
\[
f'(x) = \sin x \lim_{h\to0}\frac{\cos h-1}{h} + \cos x
\lim_{h\to0}\frac{\sin h}{h}.
\]
The previous section established $\lim_{h\to0} \sin h/h = 1$ and
$\lim_{h\to0} (\cos h-1)/h = 0$, so $f'(x) = \sin x(0) + \cos x(1)$.

\begin{theorem}
$\dfrac{d}{dx}\sin x = \cos x$.
\end{theorem}

\begin{proof}
The computation above establishes the result directly from the
definition of the derivative.
\end{proof}

The result is remarkable because the derivative of sine is another
trigonometric function rather than an entirely different kind of
expression.

\section{The Derivative of Cosine}

The same calculation applies to $f(x) = \cos x$. By definition,
\[
f'(x) = \lim_{h\to0} \frac{\cos(x+h)-\cos x}{h}.
\]
Using the cosine sum formula, $\cos(x+h) = \cos x\cos h - \sin x\sin h$,
so
\[
f'(x) = \lim_{h\to0} \frac{\cos x\cos h - \sin x\sin h - \cos x}{h} =
\lim_{h\to0}\left[\cos x\,\frac{\cos h-1}{h} - \sin x\,\frac{\sin
h}{h}\right].
\]
Applying the two fundamental limits, $f'(x) = \cos x(0) - \sin x(1) =
-\sin x$.

\begin{theorem}
$\dfrac{d}{dx}\cos x = -\sin x$.
\end{theorem}

\begin{proof}
The result follows directly from the derivative definition and the
trigonometric sum formulas, exactly as in the derivative of sine.
\end{proof}

Together, the two derivative formulas form one of the most important
relationships in mathematics: differentiation transforms sine into
cosine and cosine into negative sine. A second differentiation yields
$\dfrac{d^2}{dx^2}\sin x = -\sin x$ and $\dfrac{d^2}{dx^2}\cos x =
-\cos x$, a property that will later explain why sine and cosine
naturally describe oscillations, waves, and periodic motion.

\begin{algorithmproc}{How to Differentiate an Expression Involving Sine and Cosine}
Apply the rules $\dfrac{d}{dx}\sin x = \cos x$ and $\dfrac{d}{dx}\cos x =
-\sin x$ directly whenever the argument is simply $x$. If the angle
appearing inside sine or cosine is itself a function of $x$, differentiate
using these rules and then multiply by the derivative of that inner
function according to the chain rule.
\end{algorithmproc}

\begin{example}
Differentiate $f(x) = \sin(2x) + 3\cos x$. Using the derivative rules
together with the chain rule, $f'(x) = 2\cos(2x) - 3\sin x$.
\end{example}

\begin{practicenow}
Differentiate $f(x) = 4\sin(3x) - 2\cos(2x)$.
\end{practicenow}

\section{Circular Motion Revisited}

Chapter~5 studied circular motion using the parametric equations $x(t) =
\cos t$, $y(t) = \sin t$, introduced geometrically as the coordinates of
a point moving around the unit circle at constant angular speed. The
derivative formulas of this chapter now allow that motion to be analyzed
quantitatively. Differentiating the coordinate functions gives $x'(t) =
-\sin t$ and $y'(t) = \cos t$, so the velocity vector is
\[
\mathbf v(t) = \begin{pmatrix} x'(t) \\ y'(t) \end{pmatrix} =
\begin{pmatrix} -\sin t \\ \cos t \end{pmatrix},
\]
while the position vector is $\mathbf r(t) = \begin{pmatrix} \cos t \\
\sin t \end{pmatrix}$. Comparing the two reveals that the velocity
vector is obtained from the position vector by a rotation through
$90^\circ$:
\[
\begin{pmatrix} -\sin t \\ \cos t \end{pmatrix} = \begin{pmatrix} 0 & -1
\\ 1 & 0 \end{pmatrix}\begin{pmatrix} \cos t \\ \sin t \end{pmatrix},
\]
and Chapter~8 identified this matrix as $R(90^\circ)$, the quarter-turn
rotation.

\begin{theorem}
For motion on the unit circle, $\mathbf r(t)\cdot\mathbf v(t) = 0$;
consequently, the velocity vector is tangent to the circle at every
point.
\end{theorem}

\begin{proof}
Compute the dot product: $\mathbf r(t)\cdot\mathbf v(t) = (\cos
t)(-\sin t) + (\sin t)(\cos t) = -\cos t\sin t + \sin t\cos t = 0$, so
the vectors are perpendicular.
\end{proof}

This result provides a satisfying closure to Chapter~5. The trigonometric
functions that originally described circular motion now explain the
geometry of that motion directly: the velocity is always tangent to the
circle, and differentiation acts as a quarter-turn rotation of the
position vector.

\section{Why Radians Matter}

Throughout this chapter, the limits $\lim_{\theta\to0} \sin\theta/\theta
= 1$ and $\lim_{\theta\to0} \tan\theta/\theta = 1$ have played a central
role, and these formulas are true only when angles are measured in
radians. To see why, suppose angle measure is expressed in degrees.
Since $180^\circ = \pi$ radians, an angle of $h$ degrees corresponds to
$\pi h/180$ radians, so
\[
\lim_{h\to0} \frac{\sin h^\circ}{h} = \lim_{h\to0}
\frac{\sin(\pi h/180)}{\pi h/180}\cdot\frac{\pi}{180}.
\]
The first factor approaches $1$, leaving $\lim_{h\to0} \sin h^\circ/h =
\pi/180$: instead of obtaining $1$, an additional constant appears, and
consequently $\dfrac{d}{dx}\sin x = \dfrac{\pi}{180}\cos x$ when $x$ is
measured in degrees. The derivative formula becomes cluttered with a
conversion factor. By contrast, when angles are measured in radians,
$\dfrac{d}{dx}\sin x = \cos x$ with no additional constant. Radians are
therefore not merely a convenient convention; they are the angular unit
naturally adapted to the geometry of circles and the calculus of
trigonometric functions, confirming the theme introduced in Chapter~4
that radians are the unit in which the fundamental structures of
trigonometry take their simplest form.

\section{Applications}

The derivative relationships $\dfrac{d}{dx}\sin x = \cos x$ and
$\dfrac{d}{dx}\cos x = -\sin x$ appear throughout mathematics, science,
and engineering. In simple harmonic motion, the position of an
oscillating object is often modeled by a sine or cosine function, while
the derivative describes its velocity. In acoustics, vibrating strings
and sound waves are represented using trigonometric functions whose
derivatives describe rates of change in pressure and displacement. In
electrical engineering, alternating-current signals are modeled by
sinusoids, with derivatives corresponding to changing voltages and
currents. In astronomy, periodic motions such as rotations and orbital
approximations are frequently expressed using trigonometric functions
and their derivatives. The reason these applications are so widespread
is now visible: differentiation preserves the trigonometric family, and
the derivative of a sine function is another trigonometric function, as
is the derivative of a cosine function.

\section{Looking Ahead}

This chapter marks a transition in the book. The earlier chapters
studied finite transformations: rotations, similarities, amplitwists,
and complex multiplication. This chapter introduced infinitesimal change
through limits and derivatives. A different kind of infinite process
remains to be studied: instead of examining what happens under
infinitely small changes, the next chapter investigates what happens
when infinitely many quantities are accumulated together. The central
objects will be sequences and series, which provide the mathematical
foundation for approximation, convergence, and many of the analytical
techniques used throughout modern mathematics.

\section{Exercises}

\subsection*{Small-Angle Limits}
\begin{exercise}
Compute $\sin\theta/\theta$ for $\theta = 0.5$, $0.1$, and $0.01$
radians.
\end{exercise}

\begin{exercise}
Compute $\tan\theta/\theta$ for $\theta = 0.5$, $0.1$, and $0.01$
radians.
\end{exercise}

\begin{exercise}
Explain numerically why $\sin\theta \approx \theta$ for small angles.
\end{exercise}

\begin{exercise}
Reproduce the geometric argument establishing $\sin\theta < \theta <
\tan\theta$.
\end{exercise}

\begin{exercise}
Use the Squeeze Theorem to prove $\lim_{\theta\to0} \sin\theta/\theta =
1$.
\end{exercise}

\subsection*{Differentiation}
\begin{exercise}
Differentiate $f(x) = \sin x$.
\end{exercise}

\begin{exercise}
Differentiate $f(x) = \cos x$.
\end{exercise}

\begin{exercise}
Differentiate $f(x) = 3\sin x - 2\cos x$.
\end{exercise}

\begin{exercise}
Differentiate $f(x) = \sin(2x)$.
\end{exercise}

\begin{exercise}
Differentiate $f(x) = 4\cos(3x)$.
\end{exercise}

\begin{exercise}
Differentiate $f(x) = \sin(5x) + \cos(2x)$.
\end{exercise}

\begin{exercise}
Differentiate $f(x) = 2\sin(4x) - 3\cos(5x)$.
\end{exercise}

\subsection*{Circular Motion}
\begin{exercise}
For $x(t) = \cos t$, $y(t) = \sin t$, find the velocity vector.
\end{exercise}

\begin{exercise}
Show that the velocity vector is perpendicular to the position vector.
\end{exercise}

\begin{exercise}
Verify that the velocity vector is tangent to the unit circle.
\end{exercise}

\begin{exercise}
Find the velocity vector at $t = \pi/2$.
\end{exercise}

\begin{exercise}
Find the velocity vector at $t = \pi$.
\end{exercise}

\subsection*{Challenge Problems}
\begin{exercise}
Prove $\lim_{\theta\to0} \tan\theta/\theta = 1$.
\end{exercise}

\begin{exercise}
Show that $\dfrac{d^2}{dx^2}\sin x = -\sin x$.
\end{exercise}

\begin{exercise}
Show that $\dfrac{d^2}{dx^2}\cos x = -\cos x$.
\end{exercise}

\begin{exercise}
Explain why the equation $y'' = -y$ is naturally solved by sine and
cosine functions.
\end{exercise}

\begin{exercise}
Represent the pair $(\cos x, \sin x)$ as a vector, and show that
differentiation acts as a quarter-turn rotation on it.
\end{exercise}

\begin{exercise}
Compare the action of differentiation on $(\cos x, \sin x)$ with
multiplication by $i$ in Chapter~16. Explain the similarity.
\end{exercise}
